\IfFileExists{stacks-project.cls}{%
\documentclass{stacks-project}
}{%
\documentclass{amsart}
}

% For dealing with references we use the comment environment
\usepackage{verbatim}
\newenvironment{reference}{\comment}{\endcomment}
%\newenvironment{reference}{}{}
\newenvironment{slogan}{\comment}{\endcomment}
\newenvironment{history}{\comment}{\endcomment}

% For commutative diagrams we use Xy-pic
\usepackage[all]{xy}

% We use 2cell for 2-commutative diagrams.
\xyoption{2cell}
\UseAllTwocells

% We use multicol for the list of chapters between chapters
\usepackage{multicol}

% This is generally recommended for better output
\usepackage{lmodern}
\usepackage[T1]{fontenc}

% For cross-file-references
\usepackage{xr-hyper}

% Package for hypertext links:
\usepackage{hyperref}

% For any local file, say "hello.tex" you want to link to please
% use \externaldocument[hello-]{hello}
\externaldocument[introduction-]{introduction}
\externaldocument[conventions-]{conventions}
\externaldocument[sets-]{sets}
\externaldocument[categories-]{categories}
\externaldocument[topology-]{topology}
\externaldocument[sheaves-]{sheaves}
\externaldocument[sites-]{sites}
\externaldocument[stacks-]{stacks}
\externaldocument[fields-]{fields}
\externaldocument[algebra-]{algebra}
\externaldocument[brauer-]{brauer}
\externaldocument[homology-]{homology}
\externaldocument[derived-]{derived}
\externaldocument[simplicial-]{simplicial}
\externaldocument[more-algebra-]{more-algebra}
\externaldocument[smoothing-]{smoothing}
\externaldocument[modules-]{modules}
\externaldocument[sites-modules-]{sites-modules}
\externaldocument[injectives-]{injectives}
\externaldocument[cohomology-]{cohomology}
\externaldocument[sites-cohomology-]{sites-cohomology}
\externaldocument[dga-]{dga}
\externaldocument[dpa-]{dpa}
\externaldocument[sdga-]{sdga}
\externaldocument[hypercovering-]{hypercovering}
\externaldocument[schemes-]{schemes}
\externaldocument[constructions-]{constructions}
\externaldocument[properties-]{properties}
\externaldocument[morphisms-]{morphisms}
\externaldocument[coherent-]{coherent}
\externaldocument[divisors-]{divisors}
\externaldocument[limits-]{limits}
\externaldocument[varieties-]{varieties}
\externaldocument[topologies-]{topologies}
\externaldocument[descent-]{descent}
\externaldocument[perfect-]{perfect}
\externaldocument[more-morphisms-]{more-morphisms}
\externaldocument[flat-]{flat}
\externaldocument[groupoids-]{groupoids}
\externaldocument[more-groupoids-]{more-groupoids}
\externaldocument[etale-]{etale}
\externaldocument[chow-]{chow}
\externaldocument[intersection-]{intersection}
\externaldocument[pic-]{pic}
\externaldocument[weil-]{weil}
\externaldocument[adequate-]{adequate}
\externaldocument[dualizing-]{dualizing}
\externaldocument[duality-]{duality}
\externaldocument[discriminant-]{discriminant}
\externaldocument[derham-]{derham}
\externaldocument[local-cohomology-]{local-cohomology}
\externaldocument[algebraization-]{algebraization}
\externaldocument[curves-]{curves}
\externaldocument[resolve-]{resolve}
\externaldocument[models-]{models}
\externaldocument[functors-]{functors}
\externaldocument[equiv-]{equiv}
\externaldocument[pione-]{pione}
\externaldocument[etale-cohomology-]{etale-cohomology}
\externaldocument[proetale-]{proetale}
\externaldocument[relative-cycles-]{relative-cycles}
\externaldocument[more-etale-]{more-etale}
\externaldocument[trace-]{trace}
\externaldocument[crystalline-]{crystalline}
\externaldocument[spaces-]{spaces}
\externaldocument[spaces-properties-]{spaces-properties}
\externaldocument[spaces-morphisms-]{spaces-morphisms}
\externaldocument[decent-spaces-]{decent-spaces}
\externaldocument[spaces-cohomology-]{spaces-cohomology}
\externaldocument[spaces-limits-]{spaces-limits}
\externaldocument[spaces-divisors-]{spaces-divisors}
\externaldocument[spaces-over-fields-]{spaces-over-fields}
\externaldocument[spaces-topologies-]{spaces-topologies}
\externaldocument[spaces-descent-]{spaces-descent}
\externaldocument[spaces-perfect-]{spaces-perfect}
\externaldocument[spaces-more-morphisms-]{spaces-more-morphisms}
\externaldocument[spaces-flat-]{spaces-flat}
\externaldocument[spaces-groupoids-]{spaces-groupoids}
\externaldocument[spaces-more-groupoids-]{spaces-more-groupoids}
\externaldocument[bootstrap-]{bootstrap}
\externaldocument[spaces-pushouts-]{spaces-pushouts}
\externaldocument[spaces-chow-]{spaces-chow}
\externaldocument[groupoids-quotients-]{groupoids-quotients}
\externaldocument[spaces-more-cohomology-]{spaces-more-cohomology}
\externaldocument[spaces-simplicial-]{spaces-simplicial}
\externaldocument[spaces-duality-]{spaces-duality}
\externaldocument[formal-spaces-]{formal-spaces}
\externaldocument[restricted-]{restricted}
\externaldocument[spaces-resolve-]{spaces-resolve}
\externaldocument[formal-defos-]{formal-defos}
\externaldocument[defos-]{defos}
\externaldocument[cotangent-]{cotangent}
\externaldocument[examples-defos-]{examples-defos}
\externaldocument[algebraic-]{algebraic}
\externaldocument[examples-stacks-]{examples-stacks}
\externaldocument[stacks-sheaves-]{stacks-sheaves}
\externaldocument[criteria-]{criteria}
\externaldocument[artin-]{artin}
\externaldocument[quot-]{quot}
\externaldocument[stacks-properties-]{stacks-properties}
\externaldocument[stacks-morphisms-]{stacks-morphisms}
\externaldocument[stacks-limits-]{stacks-limits}
\externaldocument[stacks-cohomology-]{stacks-cohomology}
\externaldocument[stacks-perfect-]{stacks-perfect}
\externaldocument[stacks-introduction-]{stacks-introduction}
\externaldocument[stacks-more-morphisms-]{stacks-more-morphisms}
\externaldocument[stacks-geometry-]{stacks-geometry}
\externaldocument[moduli-]{moduli}
\externaldocument[moduli-curves-]{moduli-curves}
\externaldocument[examples-]{examples}
\externaldocument[exercises-]{exercises}
\externaldocument[guide-]{guide}
\externaldocument[desirables-]{desirables}
\externaldocument[coding-]{coding}
\externaldocument[obsolete-]{obsolete}
\externaldocument[fdl-]{fdl}
\externaldocument[index-]{index}

% Theorem environments.
%
\theoremstyle{plain}
\newtheorem{theorem}[subsection]{Theorem}
\newtheorem{proposition}[subsection]{Proposition}
\newtheorem{lemma}[subsection]{Lemma}

\theoremstyle{definition}
\newtheorem{definition}[subsection]{Definition}
\newtheorem{example}[subsection]{Example}
\newtheorem{exercise}[subsection]{Exercise}
\newtheorem{situation}[subsection]{Situation}

\theoremstyle{remark}
\newtheorem{remark}[subsection]{Remark}
\newtheorem{remarks}[subsection]{Remarks}

\numberwithin{equation}{subsection}

% Macros
%
\def\lim{\mathop{\mathrm{lim}}\nolimits}
\def\colim{\mathop{\mathrm{colim}}\nolimits}
\def\Spec{\mathop{\mathrm{Spec}}}
\def\Hom{\mathop{\mathrm{Hom}}\nolimits}
\def\Ext{\mathop{\mathrm{Ext}}\nolimits}
\def\SheafHom{\mathop{\mathcal{H}\!\mathit{om}}\nolimits}
\def\SheafExt{\mathop{\mathcal{E}\!\mathit{xt}}\nolimits}
\def\Sch{\mathit{Sch}}
\def\Mor{\mathop{\mathrm{Mor}}\nolimits}
\def\Ob{\mathop{\mathrm{Ob}}\nolimits}
\def\Sh{\mathop{\mathit{Sh}}\nolimits}
\def\NL{\mathop{N\!L}\nolimits}
\def\CH{\mathop{\mathrm{CH}}\nolimits}
\def\proetale{{pro\text{-}\acute{e}tale}}
\def\etale{{\acute{e}tale}}
\def\QCoh{\mathit{QCoh}}
\def\Ker{\mathop{\mathrm{Ker}}}
\def\Im{\mathop{\mathrm{Im}}}
\def\Coker{\mathop{\mathrm{Coker}}}
\def\Coim{\mathop{\mathrm{Coim}}}

% Boxtimes
%
\DeclareMathSymbol{\boxtimes}{\mathbin}{AMSa}{"02}

%
% Macros for moduli stacks/spaces
%
\def\QCohstack{\mathcal{QC}\!\mathit{oh}}
\def\Cohstack{\mathcal{C}\!\mathit{oh}}
\def\Spacesstack{\mathcal{S}\!\mathit{paces}}
\def\Quotfunctor{\mathrm{Quot}}
\def\Hilbfunctor{\mathrm{Hilb}}
\def\Curvesstack{\mathcal{C}\!\mathit{urves}}
\def\Polarizedstack{\mathcal{P}\!\mathit{olarized}}
\def\Complexesstack{\mathcal{C}\!\mathit{omplexes}}
% \Pic is the operator that assigns to X its picard group, usage \Pic(X)
% \Picardstack_{X/B} denotes the Picard stack of X over B
% \Picardfunctor_{X/B} denotes the Picard functor of X over B
\def\Pic{\mathop{\mathrm{Pic}}\nolimits}
\def\Picardstack{\mathcal{P}\!\mathit{ic}}
\def\Picardfunctor{\mathrm{Pic}}
\def\Deformationcategory{\mathcal{D}\!\mathit{ef}}
\usepackage{cleveref}
\crefname{originallemma}{original lemma}{original lemmas}
\Crefname{originallemma}{Original lemma}{Original lemmas}
\makeatletter
\newlabel{sites-cohomology-lemma-torsion}{{0DD7}{source}{}{sourcepointer-1}{}}
\newlabel{sites-cohomology-lemma-torsion@cref}{{[originallemma][0][]0DD7}{[1][0][]source}{}{}{}}
\makeatother
\crefname{originallemma}{original lemma}{original lemmas}
\Crefname{originallemma}{Original lemma}{Original lemmas}
\makeatletter
\newlabel{derived-lemma-leray-acyclicity}{{015E}{source}{}{sourcepointer-2}{}}
\newlabel{derived-lemma-leray-acyclicity@cref}{{[originallemma][0][]015E}{[1][0][]source}{}{}{}}
\makeatother
\crefname{originallemma}{original lemma}{original lemmas}
\Crefname{originallemma}{Original lemma}{Original lemmas}
\makeatletter
\newlabel{sites-cohomology-lemma-relative-Leray}{{0734}{source}{}{sourcepointer-3}{}}
\newlabel{sites-cohomology-lemma-relative-Leray@cref}{{[originallemma][0][]0734}{[1][0][]source}{}{}{}}
\makeatother
\crefname{originallemma}{original lemma}{original lemmas}
\Crefname{originallemma}{Original lemma}{Original lemmas}
\makeatletter
\newlabel{etale-cohomology-lemma-V-C-all-n-etale-ph}{{0DE4}{source}{}{sourcepointer-4}{}}
\newlabel{etale-cohomology-lemma-V-C-all-n-etale-ph@cref}{{[originallemma][0][]0DE4}{[1][0][]source}{}{}{}}
\makeatother
\crefname{originallemma}{original lemma}{original lemmas}
\Crefname{originallemma}{Original lemma}{Original lemmas}
\makeatletter
\newlabel{sites-cohomology-lemma-higher-direct-images}{{072W}{source}{}{sourcepointer-5}{}}
\newlabel{sites-cohomology-lemma-higher-direct-images@cref}{{[originallemma][0][]072W}{[1][0][]source}{}{}{}}
\makeatother
\crefname{originallemma}{original lemma}{original lemmas}
\Crefname{originallemma}{Original lemma}{Original lemmas}
\makeatletter
\newlabel{sites-cohomology-lemma-cohomology-of-open}{{03F3}{source}{}{sourcepointer-6}{}}
\newlabel{sites-cohomology-lemma-cohomology-of-open@cref}{{[originallemma][0][]03F3}{[1][0][]source}{}{}{}}
\makeatother
\crefname{originallemma}{original lemma}{original lemmas}
\Crefname{originallemma}{Original lemma}{Original lemmas}
\makeatletter
\newlabel{etale-cohomology-lemma-describe-pullback-pi-ph}{{0DDW}{source}{}{sourcepointer-7}{}}
\newlabel{etale-cohomology-lemma-describe-pullback-pi-ph@cref}{{[originallemma][0][]0DDW}{[1][0][]source}{}{}{}}
\makeatother
\crefname{originallemma}{original lemma}{original lemmas}
\Crefname{originallemma}{Original lemma}{Original lemmas}
\makeatletter
\newlabel{sites-definition-special-cocontinuous-functor}{{03CG}{source}{}{sourcepointer-8}{}}
\newlabel{sites-definition-special-cocontinuous-functor@cref}{{[originallemma][0][]03CG}{[1][0][]source}{}{}{}}
\makeatother
\crefname{originallemma}{original lemma}{original lemmas}
\Crefname{originallemma}{Original lemma}{Original lemmas}
\makeatletter
\newlabel{sites-lemma-equivalence}{{03A0}{source}{}{sourcepointer-9}{}}
\newlabel{sites-lemma-equivalence@cref}{{[originallemma][0][]03A0}{[1][0][]source}{}{}{}}
\makeatother
\crefname{originallemma}{original lemma}{original lemmas}
\Crefname{originallemma}{Original lemma}{Original lemmas}
\makeatletter
\newlabel{spaces-topologies-lemma-characterize-sheaf}{{0DFM}{source}{}{sourcepointer-10}{}}
\newlabel{spaces-topologies-lemma-characterize-sheaf@cref}{{[originallemma][0][]0DFM}{[1][0][]source}{}{}{}}
\makeatother
\crefname{originallemma}{original lemma}{original lemmas}
\Crefname{originallemma}{Original lemma}{Original lemmas}
\makeatletter
\newlabel{spaces-topologies-lemma-zariski-etale-smooth-syntomic-fppf-ph}{{0DFH}{source}{}{sourcepointer-11}{}}
\newlabel{spaces-topologies-lemma-zariski-etale-smooth-syntomic-fppf-ph@cref}{{[originallemma][0][]0DFH}{[1][0][]source}{}{}{}}
\makeatother
\crefname{originallemma}{original lemma}{original lemmas}
\Crefname{originallemma}{Original lemma}{Original lemmas}
\makeatletter
\newlabel{etale-cohomology-lemma-describe-pullback}{{09XN}{source}{}{sourcepointer-12}{}}
\newlabel{etale-cohomology-lemma-describe-pullback@cref}{{[originallemma][0][]09XN}{[1][0][]source}{}{}{}}
\makeatother
\crefname{originallemma}{original lemma}{original lemmas}
\Crefname{originallemma}{Original lemma}{Original lemmas}
\makeatletter
\newlabel{spaces-properties-lemma-stalk-pullback}{{04K2}{source}{}{sourcepointer-13}{}}
\newlabel{spaces-properties-lemma-stalk-pullback@cref}{{[originallemma][0][]04K2}{[1][0][]source}{}{}{}}
\makeatother
\crefname{originallemma}{original lemma}{original lemmas}
\Crefname{originallemma}{Original lemma}{Original lemmas}
\makeatletter
\newlabel{spaces-cohomology-lemma-weak-chow}{{089J}{source}{}{sourcepointer-14}{}}
\newlabel{spaces-cohomology-lemma-weak-chow@cref}{{[originallemma][0][]089J}{[1][0][]source}{}{}{}}
\makeatother
\crefname{originallemma}{original lemma}{original lemmas}
\Crefname{originallemma}{Original lemma}{Original lemmas}
\makeatletter
\newlabel{topologies-lemma-surjective-proper-ph}{{0DES}{source}{}{sourcepointer-15}{}}
\newlabel{topologies-lemma-surjective-proper-ph@cref}{{[originallemma][0][]0DES}{[1][0][]source}{}{}{}}
\makeatother
\title{Cohomological descent for the ph topology: pi\_X\textasciicircum{}\{-1\}K = R epsilon\_\{X,*\} a\_X\textasciicircum{}\{-1\}K}
\author{The Stacks Project Authors}
\date{stacks-project:spaces-more-cohomology.tex:local-snapshot-20261012; GNU FDL 1.2 or later}
\begin{document}
\maketitle
\noindent\textit{Editorial note.} The counted statement is the author's own lemma environment labelled lemma-cohomological-descent-etale-ph (source lines 1962-1973); the proof environment beginning on the line after the following blank line (source line 1975) is its own proof and closes at source line 2005.\par\medskip
\noindent\textit{Editorial note.} The proof invokes two results of the same chapter by label, both retained as context with the generator's transitive expansion (note that lemma-comparison-ph-etale is itself a counted card of this lane on the disjoint interval 1888-1955); four cross-chapter results are each declared as a tagged source pointer.\par\medskip
\noindent\textit{Editorial note.} The author proves only the second statement and says the first is easier and proved in exactly the same manner; that remark is reproduced verbatim and the card supplies nothing in its place.\par\medskip
\noindent\textit{Editorial note.} The retained preamble is the chapter's own (sources/stacks-spaces-more-cohomology/preamble.tex, taken with its real line count), a verbatim excerpt of the native Xy-pic setup and theorem declarations; no macro was invented or redefined.\par\medskip
\noindent\textit{Editorial scope.} One original result of the Stacks Project's chapter Cohomology of Algebraic Spaces, counted once: cohomological descent for the ph topology, i.e. pi\_X\textasciicircum{}\{-1\}K = R epsilon\_\{X,*\}a\_X\textasciicircum{}\{-1\}K and K = R a\_\{X,*\}a\_X\textasciicircum{}\{-1\}K for torsion complexes.\par\medskip
\noindent\textit{Allowed standard prerequisites (editorial).}\begin{itemize}
\item Torsion complexes and their reduction to sheaves. A complex K in D\textasciicircum{}+(X\_etale) with torsion cohomology sheaves can be represented by a bounded below complex of torsion abelian sheaves (Cohomology on Sites, declared as a tagged source pointer).. It is the reduction at source lines 1977-1981.
\item Leray's acyclicity lemma. If a complex is acyclic for the higher direct images of the functor, then the higher direct images of the complex vanish in positive degrees (Derived Categories, declared as a tagged source pointer).. It finishes the complex case at source lines 1985-1988.
\item The comparison lemma for the ph topology. a\_\{X,*\}a\_X\textasciicircum{}\{-1\}F = F and R\textasciicircum{}i epsilon\_\{X,*\}(a\_X\textasciicircum{}\{-1\}F) = 0 for i > 0 for torsion F (the author's lemma-comparison-ph-etale, retained as context).. It gives the identity and one vanishing at source lines 1992-1999.
\item Descent along the projection to the small etale site. R\textasciicircum{}j pi\_\{X,*\}(pi\_X\textasciicircum{}\{-1\}F) = 0 for j > 0 (the author's lemma-cohomological-descent-etale, retained as context).. It supplies the other vanishing in the Leray decomposition at source lines 2002-2003.
\item The relative Leray spectral sequence. The Leray spectral sequence for a composite functor, used for a\_X = epsilon\_X o pi\_X (Cohomology on Sites, declared as a tagged source pointer).. It combines the two vanishing statements at source lines 1995-1997.
\end{itemize}\medskip
\ifdefined\hypertarget\hypertarget{sourcepointer-1}{}\fi
\noindent\textit{Source reference.} Original lemma 0DD7: sources/stacks-chow/tags.txt, source lines 16757--16757. Cross-chapter pointer whose target carries the canonical tag 0DD7 in the project tag list; declared so the reference resolves inside the assembled package, while every step the counted argument uses remains inside the retained spans.\par\medskip
\ifdefined\hypertarget\hypertarget{sourcepointer-2}{}\fi
\noindent\textit{Source reference.} Original lemma 015E: sources/stacks-chow/tags.txt, source lines 1505--1505. Cross-chapter pointer whose target carries the canonical tag 015E in the project tag list; declared so the reference resolves inside the assembled package, while every step the counted argument uses remains inside the retained spans.\par\medskip
\ifdefined\hypertarget\hypertarget{sourcepointer-3}{}\fi
\noindent\textit{Source reference.} Original lemma 0734: sources/stacks-chow/tags.txt, source lines 9048--9048. Cross-chapter pointer whose target carries the canonical tag 0734 in the project tag list; declared so the reference resolves inside the assembled package, while every step the counted argument uses remains inside the retained spans.\par\medskip
\ifdefined\hypertarget\hypertarget{sourcepointer-4}{}\fi
\noindent\textit{Source reference.} Original lemma 0DE4: sources/stacks-chow/tags.txt, source lines 16789--16789. Cross-chapter pointer whose target carries the canonical tag 0DE4 in the project tag list; declared so the reference resolves inside the assembled package, while every step the counted argument uses remains inside the retained spans.\par\medskip
\ifdefined\hypertarget\hypertarget{sourcepointer-5}{}\fi
\noindent\textit{Source reference.} Original lemma 072W: sources/stacks-chow/tags.txt, source lines 9040--9040. Cross-chapter pointer whose target carries the canonical tag 072W in the project tag list; declared so the reference resolves inside the assembled package, while every step the counted argument uses remains inside the retained spans.\par\medskip
\ifdefined\hypertarget\hypertarget{sourcepointer-6}{}\fi
\noindent\textit{Source reference.} Original lemma 03F3: sources/stacks-chow/tags.txt, source lines 4478--4478. Cross-chapter pointer whose target carries the canonical tag 03F3 in the project tag list; declared so the reference resolves inside the assembled package, while every step the counted argument uses remains inside the retained spans.\par\medskip
\ifdefined\hypertarget\hypertarget{sourcepointer-7}{}\fi
\noindent\textit{Source reference.} Original lemma 0DDW: sources/stacks-chow/tags.txt, source lines 16781--16781. Cross-chapter pointer whose target carries the canonical tag 0DDW in the project tag list; declared so the reference resolves inside the assembled package, while every step the counted argument uses remains inside the retained spans.\par\medskip
\ifdefined\hypertarget\hypertarget{sourcepointer-8}{}\fi
\noindent\textit{Source reference.} Original lemma 03CG: sources/stacks-chow/tags.txt, source lines 4383--4383. Cross-chapter pointer whose target carries the canonical tag 03CG in the project tag list; declared so the reference resolves inside the assembled package, while every step the counted argument uses remains inside the retained spans.\par\medskip
\ifdefined\hypertarget\hypertarget{sourcepointer-9}{}\fi
\noindent\textit{Source reference.} Original lemma 03A0: sources/stacks-chow/tags.txt, source lines 4291--4291. Cross-chapter pointer whose target carries the canonical tag 03A0 in the project tag list; declared so the reference resolves inside the assembled package, while every step the counted argument uses remains inside the retained spans.\par\medskip
\ifdefined\hypertarget\hypertarget{sourcepointer-10}{}\fi
\noindent\textit{Source reference.} Original lemma 0DFM: sources/stacks-chow/tags.txt, source lines 16842--16842. Cross-chapter pointer whose target carries the canonical tag 0DFM in the project tag list; declared so the reference resolves inside the assembled package, while every step the counted argument uses remains inside the retained spans.\par\medskip
\ifdefined\hypertarget\hypertarget{sourcepointer-11}{}\fi
\noindent\textit{Source reference.} Original lemma 0DFH: sources/stacks-chow/tags.txt, source lines 16837--16837. Cross-chapter pointer whose target carries the canonical tag 0DFH in the project tag list; declared so the reference resolves inside the assembled package, while every step the counted argument uses remains inside the retained spans.\par\medskip
\ifdefined\hypertarget\hypertarget{sourcepointer-12}{}\fi
\noindent\textit{Source reference.} Original lemma 09XN: sources/stacks-chow/tags.txt, source lines 12536--12536. Cross-chapter pointer whose target carries the canonical tag 09XN in the project tag list; declared so the reference resolves inside the assembled package, while every step the counted argument uses remains inside the retained spans.\par\medskip
\ifdefined\hypertarget\hypertarget{sourcepointer-13}{}\fi
\noindent\textit{Source reference.} Original lemma 04K2: sources/stacks-chow/tags.txt, source lines 5958--5958. Cross-chapter pointer whose target carries the canonical tag 04K2 in the project tag list; declared so the reference resolves inside the assembled package, while every step the counted argument uses remains inside the retained spans.\par\medskip
\ifdefined\hypertarget\hypertarget{sourcepointer-14}{}\fi
\noindent\textit{Source reference.} Original lemma 089J: sources/stacks-chow/tags.txt, source lines 10499--10499. Cross-chapter pointer whose target carries the canonical tag 089J in the project tag list; declared so the reference resolves inside the assembled package, while every step the counted argument uses remains inside the retained spans.\par\medskip
\ifdefined\hypertarget\hypertarget{sourcepointer-15}{}\fi
\noindent\textit{Source reference.} Original lemma 0DES: sources/stacks-chow/tags.txt, source lines 16812--16812. Cross-chapter pointer whose target carries the canonical tag 0DES in the project tag list; declared so the reference resolves inside the assembled package, while every step the counted argument uses remains inside the retained spans.\par\medskip
\begin{lemma}
\label{lemma-comparison-ph-etale}
Let $S$ be a scheme. Let $X$ be an algebraic space over $S$.
\begin{enumerate}
\item For $\mathcal{F} \in \Sh(X_\etale)$ we have
$\epsilon_{X, *}a_X^{-1}\mathcal{F} = \pi_X^{-1}\mathcal{F}$
and $a_{X, *}a_X^{-1}\mathcal{F} = \mathcal{F}$.
\item For $\mathcal{F} \in \textit{Ab}(X_\etale)$ torsion we have
$R^i\epsilon_{X, *}(a_X^{-1}\mathcal{F}) = 0$ for $i > 0$.
\end{enumerate}
\end{lemma}

\begin{proof}
We have $a_X^{-1}\mathcal{F} = \epsilon_X^{-1} \pi_X^{-1}\mathcal{F}$.
By Lemma \ref{lemma-describe-pullback} the \'etale sheaf
$\pi_X^{-1}\mathcal{F}$ is a sheaf for the ph topology
and therefore is equal to $a_X^{-1}\mathcal{F}$ (as pulling
back by $\epsilon_X$ is given by ph sheafification).
Recall moreover that $\epsilon_{X, *}$ is the identity
on underlying presheaves.
Now part (1) is immediate from the explicit description of $\pi_X^{-1}$
in Lemma \ref{lemma-describe-pullback}.

\medskip\noindent
We will prove part (2) by reducing it to the case of schemes --
see part (1) of
\'Etale Cohomology, Lemma \ref{etale-cohomology-lemma-V-C-all-n-etale-ph}.
This will ``clearly work'' as every algebraic space is
\'etale locally a scheme. The details are given below but we urge
the reader to skip the proof.

\medskip\noindent
For an abelian sheaf $\mathcal{H}$ on $(\textit{Spaces}/X)_{ph}$ the
higher direct image $R^p\epsilon_{X, *}\mathcal{H}$ is the sheaf
associated to the presheaf $U \mapsto H^p_{ph}(U, \mathcal{H})$
on $(\textit{Spaces}/X)_\etale$. See
Cohomology on Sites, Lemma \ref{sites-cohomology-lemma-higher-direct-images}.
Since every object of $(\textit{Spaces}/X)_\etale$ has a covering
by schemes, it suffices to prove that given $U/X$ a scheme and
$\xi \in H^p_{ph}(U, a_X^{-1}\mathcal{F})$ we can find
an \'etale covering $\{U_i \to U\}$ such that $\xi$
restricts to zero on $U_i$. We have
\begin{align*}
H^p_{ph}(U, a_X^{-1}\mathcal{F})
& =
H^p((\textit{Spaces}/U)_{ph}, (a_X^{-1}\mathcal{F})|_{\textit{Spaces}/U}) \\
& =
H^p((\Sch/U)_{ph}, (a_X^{-1}\mathcal{F})|_{\Sch/U})
\end{align*}
where the second identification is
Lemma \ref{lemma-compare-cohomology-other-topologies}
and the first is a general fact about restriction
(Cohomology on Sites, Lemma \ref{sites-cohomology-lemma-cohomology-of-open}).
Looking at the first paragraph and the corresponding result in the
case of schemes (\'Etale Cohomology, Lemma
\ref{etale-cohomology-lemma-describe-pullback-pi-ph})
we conclude that the sheaf $(a_X^{-1}\mathcal{F})|_{\Sch/U}$
matches the pullback by the ``schemes version of $a_U$''.
Therefore we can find an \'etale covering
$\{U_i \to U\}$ such that our class dies in
$H^p((\Sch/U_i)_{ph}, (a_X^{-1}\mathcal{F})|_{\Sch/U_i})$
for each $i$, see
\'Etale Cohomology, Lemma \ref{etale-cohomology-lemma-V-C-all-n-etale-ph}
(the precise statement one should use here is that $V_n$ holds for all $n$
which is the statement of part (2) for the case of schemes).
Transporting back (using the same formulas as above but now for
$U_i$) we conclude $\xi$ restricts to zero over $U_i$ as desired.
\end{proof}

\begin{lemma}
\label{lemma-cohomological-descent-etale}
Let $S$ be a scheme. Let $X$ be an algebraic space over $S$.
For $K \in D(X_\etale)$ the map
$$
K \longrightarrow R\pi_{X, *}\pi_X^{-1}K
$$
is an isomorphism where
$\pi_X : \Sh((\textit{Spaces}/X)_\etale) \to \Sh(X_\etale)$ is as above.
\end{lemma}

\begin{proof}
This is true because both $\pi_X^{-1}$ and $\pi_{X, *} = i_X^{-1}$
are exact functors and the composition $\pi_{X, *} \circ \pi_X^{-1}$
is the identity functor.
\end{proof}

\begin{lemma}
\label{lemma-compare-cohomology-other-topologies}
Let $S$ be a scheme. Let $\tau \in \{\etale, fppf, ph\}$ (add more here).
The inclusion functor
$$
(\Sch/S)_\tau \longrightarrow (\textit{Spaces}/S)_\tau
$$
is a special cocontinuous functor
(Sites, Definition \ref{sites-definition-special-cocontinuous-functor})
and hence identifies topoi.
\end{lemma}

\begin{proof}
The conditions of Sites, Lemma \ref{sites-lemma-equivalence}
are immediately verified as our functor is fully faithful
and as every algebraic space has an \'etale covering by schemes.
\end{proof}

\begin{lemma}
\label{lemma-describe-pullback}
Let $S$ be a scheme. Let $X$ be an algebraic space over $S$.
Let $\mathcal{F}$ be a sheaf on $X_\etale$. Then
$\pi_X^{-1}\mathcal{F}$ is given by the rule
$$
(\pi_X^{-1}\mathcal{F})(Y) = \Gamma(Y_\etale, f_{small}^{-1}\mathcal{F})
$$
for $f : Y \to X$ in $(\textit{Spaces}/X)_\etale$.
Moreover, $\pi_Y^{-1}\mathcal{F}$ satisfies the
sheaf condition with respect to smooth, syntomic, fppf, fpqc, and ph coverings.
\end{lemma}

\begin{proof}
Since pullback is transitive and $f_{small} = \pi_X \circ i_f$
(see above) we see that
$i_f^{-1} \pi_X^{-1}\mathcal{F} = f_{small}^{-1}\mathcal{F}$.
This shows that $\pi_X^{-1}$ has the description given in the lemma.

\medskip\noindent
To prove that $\pi_X^{-1}\mathcal{F}$ is a sheaf for the ph topology
it suffices by Topologies on Spaces, Lemma
\ref{spaces-topologies-lemma-characterize-sheaf}
to show that for a surjective proper morphism
$V \to U$ of algebraic spaces over $X$ we have
$(\pi_X^{-1}\mathcal{F})(U)$ is the equalizer of the two maps
$(\pi_X^{-1}\mathcal{F})(V) \to (\pi_X^{-1}\mathcal{F})(V \times_U V)$.
This we have seen in Lemma \ref{lemma-surjective-proper}.

\medskip\noindent
The case of smooth, syntomic, fppf coverings follows from the case
of ph coverings by Topologies on Spaces, Lemma
\ref{spaces-topologies-lemma-zariski-etale-smooth-syntomic-fppf-ph}.

\medskip\noindent
Let $\mathcal{U} = \{U_i \to U\}_{i \in I}$ be an fpqc covering of algebraic
spaces over $X$. Let $s_i \in (\pi_X^{-1}\mathcal{F})(U_i)$ be sections
which agree over $U_i \times_U U_j$. We have to prove there exists a unique
$s \in (\pi_X^{-1}\mathcal{F})(U)$ restricting to $s_i$ over $U_i$.
Case I: $U$ and $U_i$ are schemes. This case follows from
\'Etale Cohomology, Lemma \ref{etale-cohomology-lemma-describe-pullback}.
Case II: $U$ is a scheme. Here we choose surjective \'etale morphisms
$T_i \to U_i$ where $T_i$ is a scheme. Then $\mathcal{T} = \{T_i \to U\}$ is an
fpqc covering by schemes and by case I the result holds for $\mathcal{T}$.
We omit the verification that this implies the result for $\mathcal{U}$.
Case III: general case. Let $W \to U$ be a surjective \'etale
morphism, where $W$ is a scheme. Then $\mathcal{W} = \{U_i \times_U W \to W\}$
is an fpqc covering (by algebraic spaces) of the scheme $W$.
By case II the result hold for $\mathcal{W}$.
We omit the verification that this implies the result for $\mathcal{U}$.
\end{proof}

\begin{lemma}
\label{lemma-surjective-proper}
Let $S$ be a scheme. Let $f : Y \to X$ be a surjective proper morphism
of algebraic spaces over $S$. Let $\mathcal{F}$ be a sheaf on $X_\etale$.
Then $\mathcal{F} \to f_*f^{-1}\mathcal{F}$ is injective with
image the equalizer of the two maps
$f_*f^{-1}\mathcal{F} \to g_*g^{-1}\mathcal{F}$ where
$g$ is the structure morphism $g : Y \times_X Y \to X$.
\end{lemma}

\begin{proof}
For any surjective morphism $f : Y \to X$ of algebraic spaces over $S$,
the map $\mathcal{F} \to f_*f^{-1}\mathcal{F}$ is injective.
Namely, if $\overline{x}$ is a geometric point of $X$, then we
choose a geometric point $\overline{y}$ of $Y$ lying over $\overline{x}$
and we consider
$$
\mathcal{F}_{\overline{x}} \to
(f_*f^{-1}\mathcal{F})_{\overline{x}} \to
(f^{-1}\mathcal{F})_{\overline{y}} = \mathcal{F}_{\overline{x}}
$$
See Properties of Spaces, Lemma \ref{spaces-properties-lemma-stalk-pullback}
for the last equality.

\medskip\noindent
The second statement is local on $X$ in the \'etale topology, hence we may
and do assume $Y$ is an affine scheme.

\medskip\noindent
Choose a surjective proper morphism $Z \to Y$ where $Z$ is a scheme, see
Cohomology of Spaces, Lemma \ref{spaces-cohomology-lemma-weak-chow}.
The result for $Z \to X$ implies the result for $Y \to X$.
Since $Z \to X$ is a surjective proper morphism of schemes
and hence a ph covering
(Topologies, Lemma \ref{topologies-lemma-surjective-proper-ph})
the result for $Z \to X$ follows from
\'Etale Cohomology, Lemma \ref{etale-cohomology-lemma-describe-pullback-pi-ph}
(in fact it is in some sense equivalent to this lemma).
\end{proof}

\begin{lemma}
\label{lemma-cohomological-descent-etale-ph}
Let $S$ be a scheme. Let $X$ be an algebraic space over $S$.
For $K \in D^+(X_\etale)$ with torsion cohomology sheaves the maps
$$
\pi_X^{-1}K \longrightarrow R\epsilon_{X, *}a_X^{-1}K
\quad\text{and}\quad
K \longrightarrow Ra_{X, *}a_X^{-1}K
$$
are isomorphisms with
$a_X : \Sh((\textit{Spaces}/X)_{ph}) \to \Sh(X_\etale)$ as above.
\end{lemma}

\begin{proof}
We only prove the second statement; the first is easier and proved in exactly
the same manner. There is a reduction to the case where
$K$ is given by a single torsion abelian sheaf. Namely, represent $K$
by a bounded below complex $\mathcal{F}^\bullet$ of torsion
abelian sheaves. This is possible by Cohomology on Sites, Lemma
\ref{sites-cohomology-lemma-torsion}.
By the case of a sheaf we see that
$\mathcal{F}^n = a_{X, *} a_X^{-1} \mathcal{F}^n$
and that the sheaves $R^qa_{X, *}a_X^{-1}\mathcal{F}^n$
are zero for $q > 0$. By Leray's acyclicity lemma
(Derived Categories, Lemma \ref{derived-lemma-leray-acyclicity})
applied to $a_X^{-1}\mathcal{F}^\bullet$
and the functor $a_{X, *}$ we conclude. From now on assume
$K = \mathcal{F}$ where $\mathcal{F}$ is a torsion abelian sheaf.

\medskip\noindent
By Lemma \ref{lemma-comparison-ph-etale} we have
$a_{X, *}a_X^{-1}\mathcal{F} = \mathcal{F}$. Thus it suffices to show that
$R^qa_{X, *}a_X^{-1}\mathcal{F} = 0$ for $q > 0$.
For this we can use $a_X = \epsilon_X \circ \pi_X$ and
the Leray spectral sequence
(Cohomology on Sites, Lemma \ref{sites-cohomology-lemma-relative-Leray}).
By Lemma \ref{lemma-comparison-ph-etale}
we have $R^i\epsilon_{X, *}(a_X^{-1}\mathcal{F}) = 0$ for $i > 0$.
We have
$\epsilon_{X, *}a_X^{-1}\mathcal{F} = \pi_X^{-1}\mathcal{F}$
and by Lemma \ref{lemma-cohomological-descent-etale} we have
$R^j\pi_{X, *}(\pi_X^{-1}\mathcal{F}) = 0$ for $j > 0$.
This concludes the proof.
\end{proof}

\end{document}
